% chapters_parte2/02_switching_avanzado.tex

\chapter{\eh{electric-plug} Switching Avanzado}

Si en la Parte 1 aprendiste que un switch aprende MACs y reenvía tramas, aquí vas a aprender
cómo los switches de nivel empresarial manejan escenarios que un switch hogareño ni soñaría:
múltiples instancias de Spanning Tree, cifrado en el cable, VLANs anidadas, aislamiento de
puertos, control de tormentas de tráfico, y autenticación de dispositivos en el puerto físico.

% ─────────────────────────────────────────────
\section{\eh{repeat-button} MSTP — Multiple Spanning Tree Protocol (802.1s)}

\begin{concepto}
\textbf{Problema con RSTP clásico:} una sola instancia STP para \emph{todas} las VLANs.
Eso significa que si el árbol bloquea un enlace, lo bloquea para \textbf{todo el tráfico},
sin importar la VLAN. Es como cerrar un carril de una autopista y afectar a todos los
vehículos, aunque la mitad fueran por destinos distintos.

\textbf{MSTP (802.1s)} mapea grupos de VLANs a instancias STP independientes, cada una
con su propio árbol y su propio Root Bridge. Resultado: diferentes VLANs pueden usar
diferentes caminos físicos de forma simultánea.
\end{concepto}

\subsection{Conceptos clave}

\begin{itemize}
    \item \textbf{MST Region:} conjunto de switches con la misma configuración MST
          (mismo nombre de región, número de revisión y mapeo VLAN$\to$instancia). Los
          switches fuera de la región se ven como un único switch virtual (boundary).
    \item \textbf{IST (Internal Spanning Tree / Instancia 0):} instancia especial que
          interopera con switches RSTP/STP clásicos en la frontera de la región.
    \item \textbf{MSTI (MST Instance):} instancias 1--64, cada una con su propio Root
          Bridge y topología independiente.
\end{itemize}

\subsection{Configuración Cisco IOS}

\begin{codigo}
\begin{lstlisting}
spanning-tree mode mst

spanning-tree mst configuration
  name CORP_REGION
  revision 1
  instance 1 vlan 10-20
  instance 2 vlan 30-40
  exit

! Root bridge para instancia 1
spanning-tree mst 1 priority 4096

! Root bridge para instancia 2
spanning-tree mst 2 priority 4096
\end{lstlisting}
\end{codigo}

\subsection{Diagrama: MSTP balanceando tráfico}

\begin{center}
\begin{tikzpicture}[
    sw/.style={draw=cyan!60!white, fill=darkcardgray, text=textlight,
               rectangle, rounded corners=4pt, minimum width=2.4cm,
               minimum height=0.8cm, font=\small\bfseries, align=center},
    fwd1/.style={draw=blue!50!white, ultra thick},
    fwd2/.style={draw=green!50!white, ultra thick},
    blk/.style={thick, dashed, draw=orange!60!white},
]
    \node[sw] (a) at (0,3)   {SW-A\\Root MST1};
    \node[sw] (b) at (4,3)   {SW-B\\Root MST2};
    \node[sw] (c) at (2,1.5) {SW-C};
    \node[sw] (d) at (0,0)   {SW-D};
    \node[sw] (e) at (4,0)   {SW-E};

    \draw[fwd1] (a) -- node[left, font=\tiny, text=blue!60!white] {VLAN 10-20} (c);
    \draw[fwd2] (b) -- node[right, font=\tiny, text=green!60!white] {VLAN 30-40} (c);
    \draw[fwd1] (c) -- (d);
    \draw[fwd2] (c) -- (e);
    \draw[blk]  (a) -- node[below left, font=\tiny, text=orange!70!white] {BLK} (d);
    \draw[blk]  (b) -- node[below right, font=\tiny, text=orange!70!white] {BLK} (e);
\end{tikzpicture}

{\small VLAN 10-20 usan el camino izquierdo; VLAN 30-40 el derecho.
Cada instancia bloquea enlaces distintos: \textbf{cero desperdicio de ancho de banda}.}
\end{center}

% ─────────────────────────────────────────────
\section{\eh{locked} MACsec --- MAC Security (802.1AE)}

\begin{concepto}
\textbf{MACsec} cifra el tráfico \emph{hop-by-hop} en Layer 2: entre dos switches
adyacentes, entre un switch y una laptop, entre dos datacenters unidos por fibra oscura.

Analogía: TLS es el sobre sellado que va de tu laptop al servidor de Google.
MACsec es el vehículo blindado que transporta ese sobre entre cada par de edificios.
TLS protege el contenido; MACsec protege el cable físico en cada tramo.
\end{concepto}

\subsection{Cómo funciona}

\begin{enumerate}
    \item \textbf{MKA (MACsec Key Agreement):} protocolo basado en 802.1X que negocia
          claves simétricas entre los peers. Corre sobre el mismo puerto que después cifrará.
    \item \textbf{SecY (Security Entity):} bloque hardware/software que cifra/descifra
          cada frame Ethernet en tiempo real.
    \item \textbf{Overhead:} 32 bytes por frame (16 bytes SecTAG + 16 bytes ICV). Latencia
          adicional: microsegundos en hardware con ASIC dedicado.
\end{enumerate}

\subsection{Formato del frame MACsec}

\begin{verbatim}
+----------+----------+------------+-----------------+----------+----------+
| Dst MAC  | Src MAC  | 0x88E5     |    SecTAG       | Payload  |   ICV    |
| 6 bytes  | 6 bytes  | (S-Tag ET) |  8-16 bytes     | cifrado  | 16 bytes |
+----------+----------+------------+-----------------+----------+----------+
                                    TCI|AN|SL|PN
\end{verbatim}

\begin{cuidado}
MACsec requiere soporte en \textbf{hardware} (ASIC). No todos los switches lo implementan.
Cisco Catalyst 9000 y Nexus 9000 lo soportan nativamente. Switches de gama baja, no.
\end{cuidado}

\subsection{TLS vs MACsec}

\begin{center}
\begin{tabular}{lll}
\toprule
\textbf{Característica} & \textbf{TLS} & \textbf{MACsec} \\
\midrule
Capa           & L7 (aplicación)   & L2 (enlace) \\
Alcance        & End-to-end        & Hop-by-hop \\
Qué cifra      & Payload TCP/UDP   & Frame Ethernet completo \\
Visible al ISP & Solo IPs          & Nada (MACs también ocultas) \\
Requiere app   & Sí                & No (transparente a la app) \\
\bottomrule
\end{tabular}
\end{center}

% ─────────────────────────────────────────────
\section{\eh{nut-and-bolt} QinQ --- Double Tagging (802.1ad)}

\begin{concepto}
\textbf{Problema:} un cliente tiene sus propias VLANs (1--4094). El carrier (ISP) también
usa VLANs internamente. Si ambos usan VLAN 100, hay colisión: el switch del carrier
no distingue ``VLAN 100 del Cliente A'' de ``VLAN 100 del Cliente B''.

\textbf{QinQ} añade un \emph{segundo tag} de VLAN. El tag interno (C-Tag, Customer Tag)
pertenece al cliente. El tag externo (S-Tag, Service Tag) pertenece al carrier.
El carrier solo ve S-Tags; los C-Tags son opacos para él.
\end{concepto}

\subsection{Formato del frame QinQ}

\begin{verbatim}
+---------+---------+----------+---------+----------+---------+
| Dst MAC | Src MAC | 0x88A8   | S-Tag   | 0x8100   | C-Tag   |
| 6 bytes | 6 bytes |(S-Tag ET)|(4 bytes)|(C-Tag ET)|(4 bytes)|
+---------+---------+----------+---------+----------+---------+
          ^-- Carrier procesa hasta aquí
                                          ^-- Cliente ve esto
\end{verbatim}

\begin{ejemplo}
Telmex te vende metro ethernet para conectar Polanco con Vallejo. Internamente usa S-Tag
500 para identificar tu tráfico. Dentro de tu red usas VLAN 10 (datos) y VLAN 20 (voz).
Los switches de Telmex no saben nada de tus VLANs internas: solo mueven frames con S-Tag 500.
\end{ejemplo}

% ─────────────────────────────────────────────
\section{\eh{house} Private VLANs}

\begin{concepto}
\textbf{Private VLANs (PVLANs)} resuelven el aislamiento \emph{dentro} de una VLAN.
Escenario: hosting con 50 VMs del Cliente A en la misma VLAN. No quieres que se vean
entre sí — solo deben alcanzar el gateway. Con VLANs normales, cualquier VM habla con
cualquier otra. PVLANs lo impiden a nivel de switch, sin necesidad de firewall adicional.
\end{concepto}

\subsection{Tipos de puertos}

\begin{center}
\begin{tabular}{lll}
\toprule
\textbf{Tipo} & \textbf{Puede hablar con} & \textbf{Caso de uso} \\
\midrule
Promiscuous  & Todos                        & Router / gateway / firewall \\
Isolated     & Solo Promiscuous             & VM aislada, servidor de cliente \\
Community    & Su comunidad + Promiscuous   & VMs del mismo tenant que sí se comunican \\
\bottomrule
\end{tabular}
\end{center}

\begin{codigo}
\begin{lstlisting}
vlan 100
  private-vlan primary
vlan 101
  private-vlan isolated
vlan 102
  private-vlan community
vlan 100
  private-vlan association 101,102

! Puerto promiscuous (router/gateway)
interface GigabitEthernet0/1
  switchport mode private-vlan promiscuous
  switchport private-vlan mapping 100 101,102

! Puerto isolated (VM aislada)
interface GigabitEthernet0/2
  switchport mode private-vlan host
  switchport private-vlan host-association 100 101
\end{lstlisting}
\end{codigo}

% ─────────────────────────────────────────────
\section{\eh{satellite-antenna} IGMP Snooping}

\begin{concepto}
Sin IGMP Snooping, un switch trata tráfico multicast como \textbf{broadcast}: lo manda
a todos los puertos del segmento. Un stream de video 4K a 20 Mbps llega a todos,
aunque solo una laptop lo haya pedido.

Con IGMP Snooping, el switch \emph{escucha} los mensajes IGMP que cruzan sus puertos
y construye una tabla de ``qué puerto quiere qué grupo multicast''. El tráfico
solo va a los puertos interesados.
\end{concepto}

\subsection{Flujo de aprendizaje}

\begin{enumerate}
    \item Router manda \textbf{IGMP General Query} a \texttt{224.0.0.1} cada 60 s.
    \item Host interesado responde con \textbf{IGMP Membership Report} (grupo específico).
    \item Switch anota: ``Puerto X quiere grupo 239.1.1.1''.
    \item Host que ya no quiere el grupo manda \textbf{IGMP Leave}.
    \item Switch elimina la entrada (o espera timeout de membresía).
\end{enumerate}

\begin{cuidado}
IGMP Snooping necesita un \textbf{IGMP Querier} en la VLAN (normalmente el router
multicast). Sin Querier, los hosts nunca envían Reports espontáneamente y el switch
no aprende → el tráfico multicast termina siendo flooded de todas formas.

Si no hay router multicast, configura \texttt{ip igmp snooping querier} en el switch.
\end{cuidado}

% ─────────────────────────────────────────────
\section{\eh{cloud-with-lightning} Storm Control}

\begin{concepto}
Un loop de red o NIC defectuoso puede generar tráfico broadcast/multicast que se multiplica
exponencialmente: un frame genera 2, esos 2 generan 4, y en segundos tienes una
\emph{broadcast storm} que satura todos los enlaces al 100\%.

\textbf{Storm Control} monitorea el porcentaje de ancho de banda consumido por
broadcast/multicast/unknown unicast. Si supera el umbral, descarta los frames en exceso
(o apaga el puerto).
\end{concepto}

\begin{codigo}
\begin{lstlisting}
interface GigabitEthernet0/1
  storm-control broadcast level 20.00
  storm-control multicast level 10.00
  storm-control action shutdown
\end{lstlisting}
\end{codigo}

\begin{center}
\begin{tabular}{ll}
\toprule
\textbf{Acción} & \textbf{Comportamiento} \\
\midrule
\texttt{drop}     & Descarta exceso, puerto sigue activo \\
\texttt{shutdown} & Apaga el puerto (err-disabled) \\
\texttt{trap}     & Genera SNMP trap al NMS \\
\bottomrule
\end{tabular}
\end{center}

\begin{itemize}
    \item \textbf{STP} previene loops $\to$ evita que la tormenta ocurra.
    \item \textbf{Storm Control} limita el daño si la tormenta \emph{ya ocurrió}.
    Son complementarios, no alternativos.
\end{itemize}

% ─────────────────────────────────────────────
\section{\eh{key} 802.1X en Switches --- Wired NAC}

\begin{concepto}
El mismo 802.1X del Wi-Fi empresarial aplica a \textbf{puertos físicos de switch}. Un
puerto en estado \emph{unauthorized} no permite tráfico de datos hasta que el dispositivo
autentica contra RADIUS.

Analogía: el puerto del switch es el torniquete de un edificio. Sin credenciales válidas
no entras, aunque físicamente hayas insertado el cable.
\end{concepto}

\subsection{MAB --- MAC Authentication Bypass}

\begin{concepto}
Impresoras, teléfonos IP y cámaras no tienen \emph{supplicant} 802.1X. No pueden
iniciar el handshake EAP.

\textbf{MAB} es el fallback: si el dispositivo no responde al EAPOL-Request inicial,
el switch toma su MAC y la envía a RADIUS como si fuera el ``usuario''. RADIUS tiene
una base de datos de MACs autorizadas.

Limitación: MACs se pueden falsificar (\emph{MAC spoofing}). MAB no es autenticación
fuerte; es ``mejor que nada'' para dispositivos sin 802.1X.
\end{concepto}

\begin{codigo}
\begin{lstlisting}
aaa new-model
aaa authentication dot1x default group radius
aaa authorization network default group radius
dot1x system-auth-control

interface GigabitEthernet0/2
  switchport mode access
  switchport access vlan 10
  authentication port-control auto
  authentication order dot1x mab
  authentication priority dot1x mab
  dot1x pae authenticator
  mab
\end{lstlisting}
\end{codigo}

\subsection{VLANs de contingencia}

\begin{itemize}
    \item \textbf{Guest VLAN:} dispositivos que no inician 802.1X caen aquí (acceso limitado).
    \item \textbf{Auth-Fail VLAN:} intentaron autenticar pero fallaron (portal de error).
    \item \textbf{Critical VLAN:} si RADIUS está caído, el puerto cae aquí en vez de
          bloquearse por completo. Evita lockout total de la red.
\end{itemize}

\begin{moderno}
\emoji{sparkles} \textbf{Cisco ISE (Identity Services Engine)} integra 802.1X + MAB +
postura de seguridad (¿el dispositivo tiene antivirus? ¿parches aplicados?) + profiling
automático (identifica si es impresora HP o iPhone) y asigna políticas por identidad,
no por puerto físico. Es el cerebro del NAC empresarial moderno.
\end{moderno}
